\newpage
\section{Filtros para estudiar subconjuntos de nodos}

\subsection{Minired del S\&P 500}

Para estudiar subconjuntos de nodos en la minired, el filtro más informativo se escoge a partir de las puntuaciones de \textit{in-degree} y \textit{out-degree} porque es una red dirigida de proveedores a compradores.

Una forma sintética de separar roles consiste en definir
\[
 g(v)=\mathrm{k_{in}}(v)-\mathrm{k_{out}}(v).
\]
Esta diferencia no reemplaza a las métricas originales, pero resume el balance entre emisión y recepción dentro del subgrafo. Su interpretación es directa:
\begin{itemize}[leftmargin=1.5em]
    \item si $g(v)\gg 0$, el nodo opera como \emph{proveedor neto};
    \item si $g(v)\ll 0$, el nodo opera como \emph{comprador neto};
    \item si $g(v)\approx 0$, el nodo no presenta un sesgo en una de las dos direcciones.
\end{itemize}

Este filtro permite separar funciones económicas. Dos nodos pueden tener el mismo número total de enlaces y, sin embargo, ocupar papeles opuestos.

\subsection{Red de \textit{Star Wars Episode II}}

En la red de coapariciones, un filtro más útil es el basado en grado, intermediación y pertenencia al núcleo narrativo (de eigenvector). El grado permite detectar personajes con muchas conexiones directas; la intermediación identifica personajes que enlazan grupos o escenas distintas; y el núcleo narrativo agrupa a los vértices que concentran la estructura de la película.

En una red no dirigida no tiene sentido separar emisor y receptor, por lo que un filtro basado en \textit{hubs} y \textit{authorities} no añade contenido conceptual nuevo. En cambio, seleccionar personajes con betweenness alta permite distinguir protagonistas que, además de ser visibles, articulan múltiples regiones del relato. Ese criterio es adecuado porque se desea diferenciar entre popularidad local y capacidad de conexión estructural.
